\section{Signed intersections and the residual wall}\label{sec:wall-geometry}
This section distinguishes three objects: the incidence over the
rank-two base, its reduced normalization, and its image in the full
multiplication failure scheme. They agree neither globally nor in all
scheme structures.

\subsection{Normalization and the sign incidence complex}
Let $B_2^\circ$ be the open set of rank-two projective Hankel forms with
two distinct support points and nonzero weights. Choosing one of the
two isotropic hyperplanes gives a finite etale double cover
$\widetilde B_2^\circ\longrightarrow B_2^\circ$. On that cover the
hyperplane bundles $H_+,H_-$ and their common radical bundle $E$ are
labelled. Denote by $\mathfrak I_2^\circ$ the restriction of
$\mathfrak I$ to this base.
\begin{theorem}[Normalization and all signed intersections over the rank-two base]
\label{thm:signed-normalization}
After this double-cover base change, the normalization of
$(\mathfrak I_2^\circ)_{\rm red}$ is the disjoint union of the smooth
relative varieties
\[
 Y_\varepsilon=\prod_{\nu=1}^L\Gr(c,H_{\varepsilon_\nu}),
 \qquad\varepsilon\in\{+,-\}^L.
\]
Descent identifies $\varepsilon$ with $-\varepsilon$, giving
$2^{L-1}$ irreducible components before the cover.

For any nonempty collection $\mathcal A$ of sign patterns let $J$ be
the set of contacts at which those patterns do not all agree. Their
scheme-theoretic intersection as reduced components over the labelled
base is
\begin{equation}\label{eq:sign-intersections}
 \bigcap_{\varepsilon\in\mathcal A}Y_\varepsilon
 =\prod_{\nu\in J}\Gr(c,E)
  \mathop{\times}\prod_{\nu\notin J}\Gr(c,H_{\varepsilon_\nu}).
\end{equation}
When nonempty it is smooth of dimension $M-b-c|J|$ and has
codimension $c|J|$ in each $Y_\varepsilon$.
These are also the local projected intersections wherever the
multiplication map has corank one and its annihilator belongs to
$B_2^\circ$.
\end{theorem}
\begin{proof}
For a split rank-two form its quadratic restriction is a product of two
linear functionals. A vector subspace is isotropic exactly when one of
these functionals vanishes identically on it: the polynomial ring on
that subspace is a domain. This identifies the reduced incidence as the
union of the displayed $Y_\varepsilon$. Each is smooth and normal.
Their disjoint union maps finitely to the union, is birational on every
component, and hence is its normalization. The deck transformation
simultaneously interchanges the two hyperplanes; its action on sign
patterns has no fixed pattern. The cover is connected, as is also seen
from the parameterization by a secant hyperplane and its uniquely
recovered partner. Thus its orbits give the stated component count.

The condition $U\subset H_+$ and $U\subset H_-$ is exactly
$U\subset E$. On any frame chart these are independent linear
containment equations; adding the ideals of the reduced components
therefore gives precisely the Grassmannians in
\eqref{eq:sign-intersections}, with their reduced scheme structures.
Replacing $\Gr(c,k-1)$ by $\Gr(c,k-2)$ lowers dimension by $c$ at
each differing contact. This proves the formula, including emptiness
if a required $c$-plane cannot fit in $E$. The corank-one incidence
isomorphism transfers the local assertion to the projected scheme.
\end{proof}
The restriction on the last statement is essential. Two signed images
can also meet through different annihilators at a higher-corank point,
or over a support collision outside $B_2^\circ$. Formula
\eqref{eq:sign-intersections} does not identify such points with a
single labelled annihilator.

\begin{proposition}[The embedded structure of a rank-two fibre]
\label{prop:rank-two-fibre}
Fix a split rank-two form and one contact with $U\subset E$.
Up to smooth local factors, its full isotropy fibre has coordinate ring
\begin{equation}\label{eq:rank-two-fibre-ring}
 R/J_c,\quad R=\C[a_1,\ldots,a_c,b_1,\ldots,b_c],\quad
 J_c=(a_i b_i,\ a_i b_j+a_j b_i\ (i<j)).
\end{equation}
Its radical is $(a_1,\ldots,a_c)\cap(b_1,\ldots,b_c)$.
For $c\ge2$ its nilradical has dimension $\binom c2$ as a
$\C$-vector space, is annihilated by all $a_i,b_i$, and squares to
zero. Its associated primes are exactly
\[
 (a_1,\ldots,a_c),\quad (b_1,\ldots,b_c),\quad
 (a_1,\ldots,a_c,b_1,\ldots,b_c).
\]
For several contacts the fibre is the product of the corresponding
isotropy fibres, with independent contact coordinates.
\end{proposition}
\begin{proof}
Choose quotient coordinates so that the rank-two quadratic form is
$2ab$. In a Grassmannian frame chart at $U\subset E$, its two quotient
coordinates on the $i$th frame vector are $a_i,b_i$. The remaining
coordinates are unconstrained radical directions. Restricting the
bilinear form to all pairs of frame vectors gives exactly $J_c$.
Its zero set is the union of the two coordinate $c$-planes, by the
factorization argument above, proving the radical assertion.

Modulo $J_c$, the mixed quadratic monomials satisfy
$a_i b_i=0$ and $a_i b_j=-a_j b_i$. Their classes therefore form
$\bigwedge^2\C^c$, with no further degree-two relations. Every mixed
cubic vanishes: for example
\[
 a_i a_k b_j=-a_i a_j b_k=a_j a_k b_i=-a_k a_i b_j,
\]
so it is zero in characteristic zero; the argument for one $a$ and two
$b$'s is identical. All mixed monomials of higher degree vanish as
well. This proves the assertions about the nilradical $N$.
In the exact sequence
\[
 0\longrightarrow N\longrightarrow R/J_c
 \longrightarrow R/\sqrt{J_c}\longrightarrow0,
\]
$N$ has only the origin as associated prime, whereas the reduced
quotient has its two minimal primes and no embedded prime.
The associated-prime containment for a short exact sequence permits
no others, and all three occur when $c\ge2$. The equations for
different contacts use disjoint variables, proving the product
presentation.
\end{proof}
This embedded prime is a statement about a fibre with its annihilator
fixed. It does not assert an embedded prime of $\mathscr D_L$.
In particular a reduced total scheme can have these nonreduced fibres.

\subsection{Ordinary double crossings at the residual wall}
For the next theorem suppose
\begin{equation}\label{eq:wall-hypotheses}
 c\ge8,\quad k\ge2c+1,\quad a=b.
\end{equation}
The parity in $L\binom c2=2k-5$ in fact forces $c\ge10$ in this
range. On the signed open set of
Proposition~\ref{prop:residual-dominance}, the residual map is square:
\[
 \mu_K':\bigoplus_\nu\Sym^2K_\nu\longrightarrow V_{2n-4},
 \qquad L\binom c2=2k-5.
\]
Let $Z_\varepsilon$ be the inverse image of its determinant divisor.
It is a divisor in the signed incidence, rather than a new arbitrary
choice of a discriminant.

\begin{theorem}[The residual double-crossing divisor]\label{thm:wall-node}
Under \eqref{eq:wall-hypotheses}, $Z_\varepsilon$ is irreducible and
reduced. It has a nonempty dense open subset on which the following
assertions hold. The original multiplication has corank one, and the
residual multiplication has corank one with a nonsingular residual
Hankel annihilator. Put $d=M-a$. The analytic germ of $\mathscr D_L$ at its projected point has
equation $uv=0$ in smooth local coordinates
$(t_1,\ldots,t_{d-1},u,v)$, and its completed local ring is
\begin{equation}\label{eq:wall-node-ring}
 \C[[t_1,\ldots,t_{d-1},u,v]]/(uv).
\end{equation}
The two branches are $\Sigma_\varepsilon$ and $\Phi_L$; no other
signed branch passes through this open subset. Their intersection has
codimension one in either branch, and codimension $a+1$ in $X_L$.
The tangent cone is the union of two distinct hyperplanes in the
$(d+1)$-dimensional tangent space. The germ is seminormal, its
normalization separates its two smooth branches, and its conductor in
\eqref{eq:wall-node-ring} is $(u,v)$.
\end{theorem}
\begin{proof}
At residual dimensions $k'=k-2$, $c'=c-1$ the expected codimension
is one and the signed codimension is $L(c-1)-3>1$.
Here $c'\ge9$ and $k'\ge2c'+1$. The already proved
Theorem~\ref{thm:all-dimension} says that the residual determinant
is an integral reduced divisor whose general annihilator $\eta$ is
nonsingular. Pullback by the smooth affine-space bundle in
Proposition~\ref{prop:residual-dominance} preserves this generic
irreducibility and reducedness; bundle trivializations as the support
points vary give the same assertion for $Z_\varepsilon$.

We check that the original map has corank one on a nonempty open part
of this divisor. A general residual tuple has $K_1$ mapping onto the
two evaluation coordinates at $x,y$. Indeed failure of this condition
places $K_1$ in one hyperplane of a pencil, and the isotropic
Grassmannian incidence count excludes that closed subset. Choose a
nonzero
\[
 k_0\in K_1\cap gV_{n-4},\qquad k_0=gk_0'.
\]
Such vectors exist because this intersection has dimension at least
$c-3>0$. Write $U_1=gK_1+\C a_1$.
The product $(gk_0)a_1=g^2k_0'a_1$ fills the missing direction in
$g^2V_{2n-4}$ exactly when $\eta(k_0'a_1)\ne0$.
Varying the lift by $a_1\mapsto a_1+gw$, $w\in V_{n-2}$,
changes this value by $H_\eta(k_0,w)$, a nonzero functional because
$H_\eta$ is nonsingular. Thus general lifts fill this direction.
The products $(gK_1)a_1$ give both first-derivative coordinates modulo
$g^2V_{2n-4}$, since $K_1$ has full evaluation rank and $a_1$ has
nonzero values at $x,y$. The product $a_1^2$ gives one additional value
coordinate. Consequently the original product span has dimension
$(p-4)+3=p-1$. All these conditions are nonempty open conditions.

Work on this corank-one open set in the incidence. The signed
incidence $S$ is smooth of dimension $d$ in the smooth ambient
$Y=\Pj(V_{2n}^*)\times X_L$. Its codimension is
\[
 (M+p-1)-d=p-1+b=e,
\]
and the full incidence is the zero scheme of $e$ restriction equations.
The normal derivative of these equations along $S$ splits into the
$Lc$ mixed and square equations involving the nonradical vector
$a_\nu$, and the $L\binom c2=p-4$ equations on
$gK_\nu\times gK_\nu$. The first block is invertible in transverse
plane coordinates, by the dual-frame proof of
Theorem~\ref{thm:residual-tangent}. The normal space to the rank-two
Hankel base is identified with $(g^2V_{2n-4})^*$: its tangent space
is the annihilator of $g^2V_{2n-4}$ modulo $\C\ell$, as is seen by
differentiating the two evaluations and their weights. In these
normal coordinates the remaining derivative block is precisely
$(\mu_K')^*$. Hence, up to a nowhere-zero factor,
\begin{equation}\label{eq:normal-determinant}
 \det(\text{normal derivative along }S)=\det\mu_K'.
\end{equation}
At a general point of $Z_\varepsilon$ this derivative has corank one
and its determinant vanishes simply along $S$.

Eliminate $e-1$ equations and their normal coordinates by the analytic
implicit function theorem. The remaining ambient germ is smooth of
dimension $d+1$, and $S$ is the smooth hypersurface $u=0$ in it.
The last equation vanishes on $S$, so it is $uG$. Its normal derivative
along $S$ is $G|_S$; by \eqref{eq:normal-determinant} this has a simple
zero along the residual divisor. Thus $dG$ has a nonzero component
tangent to $S$, and $u,v=G$ are independent analytic coordinates.
This proves \eqref{eq:wall-node-ring}, not only an equality of tangent
dimensions. The corank-one projection identifies this germ with the
one in $X_L$.

Only one sign choice is possible locally because no $U_\nu$ is
contained in the radical and the annihilator is unique. The other
branch must therefore be $\Phi_L$, by
Theorem~\ref{thm:all-dimension}. All dimension and tangent-cone
assertions follow from the displayed local ring. Its normalization is
\[
 \C[[t_1,\ldots,t_{d-1},u]]\ \oplus\
 \C[[t_1,\ldots,t_{d-1},v]],
\]
Replacing formal by convergent power series gives the analytic
normalization. In either case the original ring consists of pairs with
the same restriction at
$u=v=0$. This identifies the conductor as $(u,v)$ and proves
seminormality: an element of the total fraction ring whose square and
cube belong to the original ring lies in each normal branch ring, and
its two restrictions agree because their squares and cubes agree.
\end{proof}
The theorem describes a specified irreducible divisor and its dense
open double-crossing locus. It does not claim that every point of
$\Phi_L\cap\Sigma_\varepsilon$ has this type, or that higher-corank
and colliding-support intersections are normal crossings.

\begin{corollary}[A logarithmic conditioning law at a crossing]
\label{cor:node-tail}
At a real point of the double-crossing locus, with its two real
branches, there are real analytic coordinates
$(z_1,\ldots,z_{a-1},u,v,t)$ in $X_L$ in which
\begin{equation}\label{eq:node-singular-value}
 \gamma_{\rm mult}\asymp
 \left(\sum_{j=1}^{a-1}z_j^2+u^2v^2\right)^{1/2}.
\end{equation}
On a sufficiently small product box, a probability law whose density
is bounded above and below by positive constants satisfies
\begin{equation}\label{eq:node-tail}
 \Pr\{\gamma_{\rm mult}\le\epsilon\}
       \asymp\epsilon^a\log(1/\epsilon),\qquad \epsilon\downarrow0.
\end{equation}
The same law holds in any real loading chart submersing onto this
neighbourhood. Constants are local and depend on the chosen chart and
density bounds.
\end{corollary}
\begin{proof}
At corank one bounded invertible row and column operations compare
$\gamma_{\rm mult}$ with the norm of the Schur-complement row.
The analytic argument in Theorem~\ref{thm:wall-node}, applied in the
ambient Grassmannian chart, makes its ideal
$(z_1,\ldots,z_{a-1},uv)$. Each finite set of generators is a bounded
analytic linear combination of the other on a smaller compact chart,
so their norms are comparable. This proves
\eqref{eq:node-singular-value}.
For the upper bound integrate first over the $a-1$ normal coordinates,
whose volume is $O(\epsilon^{a-1})$. In a fixed square the area of
$|uv|\le C\epsilon$ is $O(\epsilon\log(1/\epsilon))$, by integrating
$\min(r,C\epsilon/|u|)$ in $u$. For the lower bound restrict the
normal coordinates to a ball of radius $\epsilon/2$ and take
$|uv|\le\epsilon/2$ in a smaller fixed square. The same elementary
integral is bounded below by a positive constant times
$\epsilon\log(1/\epsilon)$. Bounded density and Jacobian factors give
the assertion. A submersion adds unconstrained smooth coordinates and
does not change the exponent or logarithm.
\end{proof}
This is different from the smooth-stratum law $\epsilon^a$.
It is a conditioning statement for a specified real design law, not a
claim that the exact count-score compression experiment has Fisher
information equal to an unwhitened squared singular value. The
covariance of that score compression also depends on the design.
